\emph{Conventions.} The commutator product of a bivector $B$ with a vector $v$ is $B \times v = B \cdot v$, and $(x \wedge y) \cdot v = x\,(y \cdot v) - y\,(x \cdot v)$. With $\nO \cdot \nInf = 1$, $Q(\nO) = Q(\nInf) = 0$ and $W \perp \nO, \nInf$, the decomposition of a bivector into $\Lambda^{2}W$, $\R\,\nO \wedge \nInf$, $\nInf \wedge W$ and $\nO \wedge W$ ($6 + 1 + 4 + 4 = 15$) is direct, so the four slots of a bivector-valued one-form are unique. The laws of Definition \claimref{obj.synthesis.gtg-flat-spacetime-gauge-fields}(e) give, for the rotation gauge, $e \mapsto U e \tilde U$ and $\Omega \mapsto U\Omega\tilde U - 2(a \cdot \nabla U)\tilde U$; since $U$ fixes $\nO$ and $\nInf$, the law for $A$ in (G) acts slot by slot: $L \mapsto U L \tilde U - 2(a \cdot \nabla U)\tilde U$, $b \mapsto b$, $\tau \mapsto U \tau \tilde U$, $S \mapsto U S \tilde U$. For an affine displacement $f(x) = Mx + c$, $\underline{f} = M$, so $\hb(a; x) \mapsto M^{-1}\hb(a; f(x))$, hence $e(a; x) \mapsto e(Ma; f(x))$, and $\Omega(a; x) \mapsto \Omega(Ma; f(x))$; for $M = \lambda\,\mathrm{id}$ both are multiplied by $\lambda$.

\emph{(i), step 1 under (P): displacement weight.} By (T) and (P) a slot component at $x$ is a polynomial $P(x; j)$ in the jet $j$ of $(e, \Omega)$ at $x$. Translations ($\lambda = 1$) carry the jet at $f(x)$ to the jet at $x$ unchanged, and (G) then requires $P(x; j) = P(x + c; j)$ for all $c$: the coefficients are constant. Under $x \mapsto \lambda x$ a component of a $k$-th derivative of $e$ or $\Omega$ is multiplied by $\lambda^{1 + k}$, while a one-form slot must be multiplied by $\lambda$. A monomial with factors of derivative orders $k_{1}, \ldots, k_{n}$ is multiplied by $\lambda^{\sum(1 + k_{i})}$, and $\sum(1 + k_{i}) \geq n$; equality with $1$ forces $n = 1$, $k_{1} = 0$ (a constant term, $n = 0$, has weight $0$ and is excluded). So each slot component is a constant linear combination of the components $e^{i}(e_{\mu})$ and $\Omega^{j}(e_{\mu})$, which may a priori mix the background index $\mu$ with the slot's own index $\nu$. Covariance under every linear displacement $x \mapsto Mx$ ($A(a) \mapsto A(Ma)$, $e(a) \mapsto e(Ma)$, $\Omega(a) \mapsto \Omega(Ma)$) makes the coefficient matrix in the background indices commute with every invertible $M$, so it is a multiple of the identity (the $\mathrm{GL}(4)$-equivariant maps of the background covector index are the multiples of the identity): each slot is $M(e(a)) + N(\Omega(a))$ with constant linear maps $M$, $N$ acting on the values only. The check enumerates the weight-one monomials up to five factors and third derivatives: without $\hb$ there are exactly two, $e$ and $\Omega$ undifferentiated; once $\hb$ is admitted (weight $k - 1$ for its $k$-th derivative) there are weight-one monomials with any number of factors ($3, 6, 14, 25, 47$ for one to five), among them the Schouten type $\hb\,\hb\,(\partial\Omega)\,e$, of frame degree $-1$. This is where (P) enters: it is what turns the weight into a bound on the degree.

\emph{(i), step 1 under (O): frame transitivity.} By (T) and (O) the value $A(a)$ at $x$ is a function of the values $e(x)$, $\Omega(x)$ (translations remove any explicit dependence on $x$, as above). Write $A(a) = \mathcal{A}[e, \Omega](a)$. Covariance under the linear displacements gives $\mathcal{A}[e \circ M, \Omega \circ M](a) = \mathcal{A}[e, \Omega](Ma)$ for every invertible $M$. Taking $M = e^{-1}$ (the frame is invertible, and identifies background and physical vectors), $\mathcal{A}[e, \Omega](a) = \mathcal{A}[\mathrm{id}, \Omega \circ e^{-1}](e(a))$, so everything is determined by $\mathcal{A}_{0}(\omega) := \mathcal{A}[\mathrm{id}, \omega]$ for $\omega$ a linear map $W \to \Lambda^{2}W$. At a point where $U = 1$ and $a \cdot \nabla U$ is arbitrary, $\omega$ shifts by an arbitrary $\xi$ while the Lorentz slot must shift by $\xi$ and the other slots must not change; so the Lorentz slot of $\mathcal{A}_{0}(\omega)$ is $\omega$ plus a part independent of $\omega$, and the other slots are independent of $\omega$. What remains, $\mathcal{A}_{0}(0)$, is a linear map on $W$ in each slot, equivariant under constant rotors. So each slot is $M(e(a)) + N(\Omega(a))$ with constant linear maps, $N_{L} = \mathrm{id}$ and $N_{b} = N_{\tau} = N_{S} = 0$, which is where step 2 begins.

\emph{(i), step 2: rotation covariance.} At a point where $U = 1$ and $a \cdot \nabla U$ is arbitrary, $\xi = -2(a \cdot \nabla U)\tilde U$ ranges over all of $\Lambda^{2}W$; the slots $b$, $\tau$, $S$ must not change, so $N_{b}$, $N_{\tau}$, $N_{S}$ vanish, and $L = M_{L}(e) + N_{L}(\Omega)$ must change by exactly $\xi$, so $N_{L} = \mathrm{id}$. For constant $U$ the maps $M$ are equivariant for the Lorentz algebra $\Lambda^{2}W$ acting by the commutator product. The check computes the equivariant linear maps exactly: $W \to W$ is one-dimensional, spanned by the identity; $W \to \R$, $W \to \Lambda^{2}W$, $\Lambda^{2}W \to W$ and $\Lambda^{2}W \to \R$ are zero; $\Lambda^{2}W \to \Lambda^{2}W$ is two-dimensional (identity and duality), which does not matter here since $N_{L} = \mathrm{id}$ is already fixed. Hence $L = \Omega$, $b = 0$, $\tau = \kappa\, e$ and $S = \rho\, e$ with constants $\kappa$, $\rho$. Finally $(-\nInf \wedge \tau) \cdot \nO = \tau$ while $\Omega \cdot \nO = 0$ and $(\nO \wedge S) \cdot \nO = 0$, so (S) gives $\kappa = 1$: $A = A_{\rho}$.

\emph{(ii).} For $v = \sigma\nInf + \mu + t\,\nO$ with $\mu \in W$:
\[
A_{\rho}(a) \times v \;=\; \Omega(a) \cdot \mu + (t - \rho\sigma)\, e(a) + (e(a) \cdot \mu)(\rho\,\nO - \nInf).
\]
At $v = I_{\rho}$ ($\sigma = 1$, $\mu = 0$, $t = \rho$) every term vanishes, for every $(e, \Omega)$, so $D_{a} I_{\rho} = a \cdot \nabla I_{\rho} + A_{\rho}(a) \times I_{\rho} = 0$. Conversely, if every $A_{\rho}(a)$ annihilates $v$, the last term gives $e(a) \cdot \mu = 0$ for all $a$, hence $\mu = 0$ since $e$ is invertible, and then $(t - \rho\sigma)e(a) = 0$ gives $t = \rho\sigma$: $v \in \R I_{\rho}$. The check confirms the annihilator line at $\rho = \tfrac{3}{14}, -\tfrac12, 0$ (normalised kernels $31e_{+} - 25e_{-}$, $3e_{+} - 5e_{-}$, $e_{+} - e_{-}$), the soldering, and that a nonzero dilation slot leaves no annihilated vector. For $\rho \neq 0$, $I_{\rho} = \nInf + \tfrac{\Lambda}{6}\nO$ with $\Lambda = 6\rho$ is the frozen point of Theorem \claimref{dark-energy.frozen-point-makes-stelle-west-vector}; that theorem enters only here, to name the output. The bijection of the $A_{\rho}$ with the connections preserving $I_{\rho}$ is clause (ii) of Theorem \claimref{synthesis.frozen-point-field-space-is-the-gtg-pair-and-its-omega-equation-is-the-torsion-equation}, and is not used in this step.

\emph{(iii).} By Theorem \claimref{synthesis.gtg-tractor-gauge-field-reproduces-tractor-connection}, $\Omega_{T}(a) = \Omega(a) - \nInf \wedge \hb^{-1}(a) + \nO \wedge \Sch(\hb^{-1}(a))$ with $\Omega$ the torsion-free field, and $\Sch(A) = \tfrac12(\mathcal{R}(A) - \tfrac16 A\,\mathcal{R})$ (Fact \claimref{synthesis.imported-gtg-dictionary}(v)). It satisfies (T), (G) and (S); it violates (O) since $\Sch$ is built from the field strength, a first derivative of $\Omega$. For (P): the curvature operator is $\mathcal{R}(B) = -R(\hb(B))$ with $R$ the field strength of $\Omega$, which does not involve $\hb$, and the contractions $\partial_{a}$ act on physical vectors with the flat product; so at fixed $\Omega$, replacing $e$ by $c\,e$ ($\hb$ by $\hb/c$) multiplies $\mathcal{R}$, its contractions and $\Sch$ by $c^{-2}$ and the slot $\Sch(e(a))$ by $c^{-1}$. A polynomial in the jets of $e$ restricted to the ray $c\,e$ is a polynomial in $c$, and $c^{-1}K$ with $K \neq 0$ is not; the check exhibits $K \neq 0$ and the exact scaling at $c = 2$ and $c = -\tfrac32$. At a frame of corank one: scaling one physical component of $e$ by $c$, the check fits $c^{2}$ times each slot component exactly by a cubic in $c$ through five points and finds no $c^{0}$ term and a nonzero $c^{1}$ term ($-\tfrac{21}{22}$ in one component), a simple pole at $c = 0$, while $A_{\rho}$ is polynomial there and its annihilator is still the line of $I_{\rho}$. In the splitting of $\hb$ the slots of $\Omega_{T}$ and $A_{\rho}$ differ only in the special-conformal one, by $\nO \wedge (\Sch - \rho\,\mathrm{id})(e(a))$, and in the Lorentz one unless $\Omega$ is torsion-free; so they agree exactly when $\Omega$ is torsion-free and $\Sch = \rho\,\mathrm{id}$, which is the Einstein condition with $\Lambda = 6\rho$ by Theorem \claimref{synthesis.gtg-tractor-gauge-field-reproduces-tractor-connection}(ii). The check builds the curvature $\mathcal{R}(a \wedge b) = P(a) \wedge b + a \wedge P(b)$ for an Einstein $P = \rho\,\mathrm{id}$ and a non-Einstein diagonal $P$, recovers $\Sch = P$ and the scalar curvature $6\,\mathrm{tr}\,P$ ($\tfrac{36}{7}$ and $\tfrac{76}{35}$), and finds the normal connection equal to $A_{\rho}$ and $I_{\rho}$-parallel in the first case only, while $A_{\rho}$ keeps $I_{\rho}$ parallel in both. For a rotor field $g$ of $\Cl(2,4)$ acting by $A \mapsto gA\tilde g - 2(a \cdot \nabla g)\tilde g$, one has $D^{g \cdot A}_{a}(g I \tilde g) = g\,(D^{A}_{a} I)\,\tilde g$, so parallel tractors correspond under gauge transformations; every $A_{\rho}$ has one, so a normal connection without one is gauge equivalent to none of them.

\emph{The trap.} The forcing (i) and the conclusion (ii) use only the gauge laws of Definition \claimref{obj.synthesis.gtg-flat-spacetime-gauge-fields} (empty past cone) and the algebra of $\Cl(2,4)$; Theorem \claimref{dark-energy.frozen-point-makes-stelle-west-vector}, whose past cone contains the imported equivalence with Einstein's equations, enters only to name $\rho = \Lambda/6$ and $V$, after $I_{\rho}$ has been obtained. The normal tractor connection is never restricted: (P) excludes it, and (iii) records where the two meet.

\emph{(iv).} The first clause is (iii). For the second, the check takes $S(a)$ an independent $W$-valued one-form with sample rational values: the four values $A(a)$ annihilate no nonzero vector, so no constant tractor is parallel, while the same construction with $S = -\tfrac12 e$ annihilates exactly the line of $I_{-1/2}$. A parallel tractor $I$, constant or not, satisfies $F(a \wedge b) \times I = [D_{a}, D_{b}]I = 0$; with sample rational first jets of $(e, \Omega, S)$ at a point, the six values $F(e_{\mu} \wedge e_{\nu}) = \partial_{\mu}A(e_{\nu}) - \partial_{\nu}A(e_{\mu}) + A(e_{\mu}) \times A(e_{\nu})$ annihilate no nonzero vector (checked), so no tractor is parallel for generic first jets, while with $S = -\tfrac12 e$ and its jets they annihilate exactly the line of $I_{-1/2}$. An independent slot is a polynomial in its own components, so (P) holds; it is (T) that excludes it. For the independence of $\Omega$: if the Lorentz slot were a function of the frame alone, satisfying (T), (G) and (P) or (O), step 1 of (i) (either route, applied with $\Omega$ absent) makes it $M_{L}(e(a))$ with $M_{L}$ constant, which has no inhomogeneous term under the rotation gauge, whereas (G) requires $-2(a \cdot \nabla U)\tilde U$; so the Lorentz slot is the independent field $\Omega$ (concretely, the torsion-free field of $\hb$ needs both $\hb$ and first derivatives). That this independence decides only the torsion equation is clause (i) of Theorem \claimref{synthesis.frozen-point-field-space-is-the-gtg-pair-and-its-omega-equation-is-the-torsion-equation}. A dynamical frozen point, a Higgs field on the unreduced connections, reduces the connection only in the infinitely stiff limit (Theorem \claimref{synthesis.frozen-point-reduces-the-connection-only-when-infinitely-stiff}). That the freely varied slot returns the Schouten slot on shell is clause (ii) of Theorem \claimref{synthesis.frozen-point-integrand-gives-einstein-only-with-the-symmetric-slot-frozen}. For the remark on the parabolic gauge: a special-conformal rotor $U = 1 + \tfrac12\nO y$ ($y \in W$) fixes $X = \nO$, so $\ip{X}{U I \tilde U} = \ip{\tilde U X U}{I} = \ip{X}{I}$, and $D^{U \cdot A}(U I \tilde U) = U (D^{A} I)\tilde U$ as in (iii); a rescaling of the frame by a non-constant factor is not a gauge transformation of the Definition, and Einstein fields are not carried to Einstein fields by it, so $A_{\rho}$ is not a function of the conformal class.

\emph{(v), argued.} Constructions from $\hb$, its inverse and their derivatives are rational in the jets of $e$ with denominators powers of $\det e$ (since $\hb = \mathrm{adj}(e)/\det e$ and each derivative of $\hb$ adds a power). If such an expression $N/(\det e)^{k}$, $N$ a polynomial, stays bounded near the smooth points of the hypersurface $\det e = 0$ (the frames of corank one), then $N$ vanishes to order $k$ there; $\det$ is an irreducible polynomial and its real smooth points are Zariski dense in its zero set, so $(\det e)^{k}$ divides $N$ and the expression is a polynomial. The walls of Theorem \claimref{dark-energy.walls-free-sector} are frames of corank one, where the stored text uses the Einstein-Cartan equations because they are polynomial in $e$; requiring the tractor gauge field to be defined there gives (P). The requirement includes $\Omega$ regular at the wall. It is not the series' stance at every wall: at the Kasner wall of Proposition \claimref{dark-energy.frozen-point-datum-kasner-wall} the Levi-Civita block has a simple pole whose curvature invariants diverge, and consequence 1 of Theorem \claimref{dark-energy.frozen-point-makes-stelle-west-vector} lets the equations reach that wall as conditions on the residue. The polynomiality stored there is that of the action in the connection, $I$ and $V$, shared by the unreduced field space; it supports (P) for the assembly map only once (T) is granted. That (W) is a principle of the series is not stated in the store; it is offered as its one available reading.

\emph{(vi).} $\ip{X}{I_{\rho}} = \ip{\nO}{\nInf + \rho\,\nO} = 1$ and $A_{\rho}$ satisfies (T), (G), (S), (P) and (O). Conversely let $A$ have nondegenerate soldering and a parallel $I = \sigma\nInf + \mu + t\,\nO$ with $\ip{X}{I} = \sigma = 1$. $Q(I)$ is constant along a parallel section, write it $2\rho$. The special-conformal rotor $U = 1 + \tfrac12\nO y$ fixes $X$ and, with $y = \mu$, carries $I$ to $\nInf + t'\,\nO$, which has $Q = 2t'$, so $t' = \rho$ (checked at a sample $I$) and the transformed tractor is the constant $I_{\rho}$; gauge transformations carry parallel tractors to parallel tractors (as in (iii)), so the transformed connection annihilates $I_{\rho}$ and is valued in its stabiliser, $\Lambda^{2}W \oplus (\rho\,\nO - \nInf) \wedge W$, hence equal to $\Omega' - \nInf \wedge e' + \rho\,\nO \wedge e'$ for a unique pair with $e'$ nondegenerate (clause (ii) of Theorem \claimref{synthesis.frozen-point-field-space-is-the-gtg-pair-and-its-omega-equation-is-the-torsion-equation}), which is $A_{\rho}$ for that pair. So the two statements of the premise single out the same connections up to gauge: the result is a premise exchange. In the language of Fact \claimref{synthesis.imported-holonomy-reductions-of-cartan-geometries}(a), the parallel tractor is a holonomy reduction of the connection to the stabiliser of $I$, and the gauge in which it is the constant $I_{\rho}$ is a reduction of the gauge group to that stabiliser.

\emph{Scope.} Proved: (i) to (iv) and (vi), with the exact checks named in the entity. Argued: (v). This is a premise exchange: an equivalent reading of the field-space premise, not a derivation of it from transport. Not addressed: why there should be an action, the matter coupling, and the value of $\rho$ (the registered input). The result is transport-convenient: the tractor frame and the comparison with the series' transport are the series', and the forcing is displacement-weight counting or frame transitivity with Lorentz representation theory.
